\section{Introduction}

Reinforcement Learning (RL) is becoming one of the most popular and successful approaches for solving continuous control
problems, especially in robotics \cite{Kober2013}.
Indeed, this setting introduces several challenges such as delayed rewards, exploration-exploitation trade-offs,
high-dimensional continuous action spaces and potentially unstable training dynamics.

Another major challenge this paper tackles is sim-to-real transfer. Policies are commonly trained in simulation, where data
collection is cheap and safe, and then applied to environments whose dynamics differ from those seen during training.
This discrepancy represents the reality gap and it can significantly degrade performance \cite{Tobin2017}.
In the second half of the project, the reality gap is represented through a variation in the mass of a manipulated object:
from a source environment with a 1 kg cube to a target environment with a 5 kg cube.

The project is divided into two complementary parts. In the first part, we implement and analyze policy-gradient methods on
the Hopper continuous-control environment, observing learning dynamics, exploration behavior, critic stability, and overall
training performance \cite{Brockman2016, Todorov2012}. Specifically, we implement REINFORCE \cite{Williams1992} with and without a constant baseline, and a one-step Actor-Critic architecture
augmented with Generalized Advantage Estimation, evaluating the impact of variance reduction techniques on learning stability and convergence.

In the second part, we investigate sim-to-real transfer on the PandaPush environment \cite{gallouedec2021pandagym}, first evaluating an on-policy approach with Proximal Policy Optimization (PPO) \cite{SchulmanPPO2017}
against Soft Actor-Critic (SAC) \cite{SAC2018}, an off-policy method, together with Uniform Domain Randomization (UDR) and Automatic Domain Randomization (ADR) \cite{rubikADR2019}
strategies to address the reality gap.

\section{Related Work}
The use of RL for robotic control has been widely studied, and a broad overview of the field can be found in \cite{Kober2013}.
The policy-gradient family starts with REINFORCE \cite{Williams1992}, which however suffers from the high variance of its Monte Carlo gradient estimates.
Baselines and actor-critic methods \cite{SuttonBarto2018} were introduced to reduce this variance, and were later improved by Generalized Advantage Estimation \cite{SchulmanGAE2015} and by trust-region or clipped objectives such as TRPO and PPO \cite{SchulmanTRPO2015, SchulmanPPO2017}.
Off-policy methods like SAC \cite{SAC2018} and TD3 \cite{Fujimoto2018} reuse past experience through a replay buffer, making them more sample efficient in continuous control.
Regarding the reality gap, domain randomization \cite{Tobin2017} trains the agent over a distribution of simulator parameters rather than a single nominal value, while ADR \cite{rubikADR2019} adapts the randomization ranges during training based on the agent's performance.
Our implementation relies on Gym and MuJoCo \cite{Brockman2016, Todorov2012}, panda-gym \cite{gallouedec2021pandagym} and Stable-Baselines3 \cite{stable-baselines3}.

